\section{Related work}
\label{sec:related}

\paragraph{Expressiveness and encodings.}
Felleisen \cite{felleisen1991expressive} made ``expressive power'' of
programming languages precise through translations that are local in the
syntax. The process-calculus community developed the study of encodings
furthest. Palamidessi \cite{palamidessi2003comparing} separated the
synchronous from the asynchronous $\pi$-calculus by showing that no
encoding of the first into the second meets a small set of reasonable
criteria, which made the choice of criteria a central question, and Gorla
\cite{gorla2010unified} proposed a set of criteria that has become a
reference point: compositionality, name invariance, operational
correspondence, divergence reflection, and success sensitiveness. Peters
\cite{peters2019comparing} surveys the criteria and how they relate. Our
result is closest in spirit to the semantic criteria. The two step
simulations of \S\ref{sec:proof} are an operational correspondence, in a
strong form with named action sequences. The \mdiv{} marker of the liveness
theorem gives divergence reflection and preservation, and \mterm{} with the
$\ev{exit}$ and $\ev{fail}$ events plays the role of success. We do not
claim the structural criteria in Gorla's form. The translation is one Step
per program point and is defined from the statement's thread-local semantics,
but it is a whole-program compilation in continuation-passing style
\cite{appel1992cps} with naming conventions for Channels, and its target is
an engine with a fixed set of primitives, not a calculus with parallel
composition. The equivalence itself is trace equivalence, which is at the
coarse end of the linear-time/branching-time spectrum
\cite{vanglabbeek2001spectrum}; Milner's weak bisimilarity
\cite{milner1989communication} is finer than the claim needs
(\S\ref{sec:liveness-statement}).

\paragraph{Workflow patterns.}
The Workflow Patterns initiative catalogued recurring control-flow
behaviors \cite{vanderaalst2003patterns}, later revised to 43 patterns with
formal definitions \cite{russell2006revised}, and the catalog became the
standard instrument for comparing workflow languages and systems
\cite{wohed2009patterns}. YAWL \cite{vanderaalst2005yawl} was designed from
the patterns, and Kiepuszewski et al.\ \cite{kiepuszewski2003fundamentals}
studied the expressive power of classes of workflow languages formally.
Pattern-based evaluation and our theorem answer different questions. A
pattern rating asks how \emph{directly} a system supports one behavior, and
it is checked against the running system. Our theorem asks whether a
\emph{class} of programs is expressible at all, with an exact notion of
equivalence, and it is checked against a model. The two are complementary,
and Dex has both (\S\ref{sec:introduction}).

\paragraph{Simulation proofs and verified systems.}
The proof technique is standard. Forward and backward simulations
\cite{lynch1995simulations} and refinement mappings with stuttering
\cite{abadi1991refinement} are the foundations; the distinction between
safety and liveness that organizes our two theorems is due to Alpern and
Schneider \cite{alpern1985liveness}. Among mechanized developments,
CompCert \cite{leroy2009compcert} proves semantic preservation for a
compiler by simulation diagrams, and relies on determinism of the target
so that preservation in one direction suffices; we use determinism of each
Dex rule for a similar economy in \S\ref{sec:dextol}. seL4
\cite{klein2009sel4} proves that a kernel implementation refines its
abstract specification. For distributed systems, Verdi
\cite{wilcox2015verdi} verifies implementations under explicit network
fault models, and IronFleet \cite{hawblitzel2015ironfleet} combines
refinement with mechanized proofs of liveness. Those
projects verify an \emph{implementation} against a specification. We verify
a property of a \emph{model} of a system, and the gap between that model
and the code (\S\ref{sec:limitations}) is the gap they close. On the
engineering side, Woos et al.\ \cite{woos2016planning} report that their
proof needed 90 invariants and that keeping it up to date as invariants
changed was the main difficulty, and they propose to plan for such change.
Our gaps of \S\ref{sec:gaps} are the same phenomenon, and the
inductiveness fuzz is our answer, in the tradition of
testing a conjecture on random inputs before proving it
\cite{claessen2000quickcheck}. The development is in Lean~4
\cite{demoura2021lean4}.

\paragraph{The source language.}
\Lang{}'s synchronization constructs go back to monitors. Condition
variables with \kw{wait} and \kw{signal} are Hoare's
\cite{hoare1974monitors}; that a woken waiter must re-acquire the lock in
competition with other threads is the discipline of Mesa
\cite{lampson1980mesa}, which \Lang{} follows, with FIFO queues and without
spurious wakeups.

\paragraph{Durable execution and workflow engines.}
The academic literature on durable execution is mostly about the
proposition of our second theorem: that a program behaves under failures
as it would without them. Sagas \cite{garciamolina1987sagas} introduced
long-lived transactions built from steps with compensations. Ramalingam and
Vaswani \cite{ramalingam2013idempotence} gave a language-based account of
fault tolerance through idempotent, logged steps, with a formal calculus
that models failures and duplicate requests. Beldi
\cite{zhang2020beldi} provides exactly-once semantics and transactions for
workflows of stateful serverless functions, and Kallas et al.\
\cite{kallas2023mu2sls} run microservice applications on a serverless
platform with formal guarantees that the generated implementations are
correct. Jangda et al.\ \cite{jangda2019serverless} give an operational
semantics of serverless platforms and characterize when a simpler, naive
semantics coincides with it. These works compare a failure-prone
execution with an ideal one for the \emph{same} program, as our failure
theorem does. Our first theorem compares two \emph{different} programming
models in the ideal setting: threads on one side, Steps on the other. We
are not aware of a prior machine-checked theorem that a durable workflow
model expresses a conventional concurrent language up to trace
equivalence. Our failure theorem differs from the works above in three
ways. It is about a model written against an existing
server, rule by rule, rather than about a calculus or a system designed
for the proof. It allows zombie Workers that keep emitting effects after
the server gave them up. And it states what repeats: every extra effect
equals an earlier one, and the first occurrences of each method's effects
come in order. Mechanized frameworks for crash safety of implementations,
such as Perennial \cite{chajed2019perennial}, address the gap between
our model and the code.
